\documentclass[11pt,a4paper]{article}

\usepackage[utf8]{inputenc}
\usepackage[T1]{fontenc}
\usepackage[spanish,es-tabla]{babel}
\usepackage[a4paper,top=2.5cm,bottom=2.3cm,left=2.2cm,right=2.2cm,headheight=28pt]{geometry}
\usepackage{array}
\usepackage{tabularx}
\usepackage{booktabs}
\usepackage{colortbl}
\usepackage{xcolor}
\usepackage{enumitem}
\usepackage{fancyhdr}
\usepackage{lastpage}
\usepackage[most]{tcolorbox}
\usepackage{csquotes}
\usepackage[backend=biber,style=ieee]{biblatex}
\usepackage[hidelinks]{hyperref}

\addbibresource{referencias.bib}

% Paleta institucional UTEQ
\definecolor{uteqgreen}{HTML}{0B7A3B}
\definecolor{uteqdark}{HTML}{075C2C}
\definecolor{midblue}{HTML}{1F5C8B}
\definecolor{rowgray}{HTML}{F2F4F3}
\definecolor{boxbg}{HTML}{F7F9F8}

\newcolumntype{Y}{>{\raggedright\arraybackslash}X}
\newcommand{\thdr}[1]{\cellcolor{uteqgreen}\textcolor{white}{\textbf{#1}}}

% Encabezado y pie de pagina
\pagestyle{fancy}
\fancyhf{}
\lhead{\footnotesize UTEQ $\cdot$ Ingenier\'ia de Requerimientos (ISR-401) $\cdot$ SPA 2026--2027}
\rhead{\footnotesize GA-TL-02 -- FORM}
\lfoot{\footnotesize PFC: Plataforma de Ventas con Recomendaci\'on IA (Quevedo)}
\rfoot{\footnotesize P\'agina \thepage\ de \pageref{LastPage}}
\renewcommand{\headrulewidth}{0.5pt}
\renewcommand{\footrulewidth}{0.4pt}

\setlength{\parindent}{0pt}
\setlength{\parskip}{5pt}
\renewcommand{\arraystretch}{1.2}

\begin{document}

% ==============================================================================
% PORTADA Y DATOS INFORMATIVOS
% ==============================================================================
\begin{titlepage}
    \centering
    \vspace*{0.6cm}
    {\Large\textbf{\textcolor{uteqdark}{UNIVERSIDAD T\'ECNICA ESTATAL DE QUEVEDO}}}\\[0.25cm]
    {\large Facultad de Ciencias de la Computaci\'on}\\[0.15cm]
    {\large Carrera de Ingenier\'ia de Software}\\[0.8cm]

    \textcolor{uteqgreen}{\rule{\linewidth}{1.5pt}}\\[0.5cm]
    {\LARGE\textbf{ACTIVIDAD FORMATIVA: GA-TL-02 -- FORM}}\\[0.3cm]
    {\Large Taller de entrevista estructurada y design thinking aplicados al PFC}\\[0.45cm]
    \textcolor{uteqgreen}{\rule{\linewidth}{1.5pt}}\\[1.0cm]

    \begin{tcolorbox}[colback=boxbg,colframe=uteqgreen,boxrule=0.8pt,arc=2mm,title={\textbf{Proyecto Fin de Curso (PFC)}}]
        \centering
        \textbf{Plataforma de Ventas con M\'odulo de Recomendaci\'on basado en Inteligencia Artificial para Establecimientos de Suministros El\'ectricos en Quevedo}
    \end{tcolorbox}

    \vspace{0.8cm}

    \noindent\begin{tabularx}{\linewidth}{|p{0.32\linewidth}|Y|}
        \hline
        \thdr{Campo} & \thdr{Detalle}\\
        \hline
        Asignatura & Ingenier\'ia de Requerimientos (ISR-401) --- 4.\textsuperscript{o} Nivel, Paralelo B\\
        \hline
        \rowcolor{rowgray} Docente & Gleiston Cicer\'on Guerrero Ulloa, PhD (\texttt{gguerrero@uteq.edu.ec})\\
        \hline
        Unidad & Unidad 2. Elicitaci\'on y An\'alisis de Requisitos\\
        \hline
        \rowcolor{rowgray} Per\'iodo Acad\'emico & Segundo Per\'iodo Acad\'emico Ordinario (SPA) 2026--2027\\
        \hline
        Semana y cierre & Semana 4 (16 al 22 de noviembre de 2026) --- Cierre: 21 de noviembre de 2026, 23:59\\
        \hline
        \rowcolor{rowgray} Integrantes del equipo &
        Gilces Carranza Jos\'e Ignacio (Guion y an\'alisis de preguntas)\\
        & V\'elez Vinces Jostin Gustavo (Mapa de empat\'ia y definici\'on del problema)\\
        & Espinoza Jeanpierre Robinson (Estructura \LaTeX{}, ensayo y custodia LOPDP)\\
        \hline
    \end{tabularx}

    \vfill
    {\normalsize Quevedo, Los R\'ios, Ecuador}\\
    {\normalsize Noviembre de 2026}
\end{titlepage}

\tableofcontents
\newpage

% ==============================================================================
% SECCION 1: CONTEXTO Y CONSENTIMIENTO INFORMADO (Integrante 3)
% ==============================================================================
\section{Contexto del taller y protocolo de consentimiento informado}

\subsection{Prop\'osito y alcance del ensayo}
Este documento recoge el dise\~no preliminar y el ensayo controlado de la entrevista estructurada y de la sesi\'on de \emph{design thinking} para el Proyecto Fin de Curso (PFC), orientados a la operaci\'on real de establecimientos de suministros el\'ectricos en el cant\'on Quevedo. La combinaci\'on de entrevista estructurada y herramientas de empat\'ia permite comparar respuestas de forma consistente y detectar necesidades operativas que los usuarios no verbalizan de forma directa \cite{pohl2025,ireb2026}.

El ensayo se ejecuta bajo modalidad de juego de roles dentro del aula, tomando como parte interesada prioritaria al \textbf{Cajero / Vendedor de Mostrador (SH-01)}, dado que este rol concentra la atenci\'on directa en horas pico, la consulta de inventario y la asesor\'ia t\'ecnica sobre compatibilidad de materiales el\'ectricos.

\subsection{Modelo de consentimiento informado (LOPDP Ecuador)}
En cumplimiento del C\'odigo de \'Etica de la ACM \cite{acm2018} y de la Ley Org\'anica de Protecci\'on de Datos Personales (LOPDP) del Ecuador \cite{lopdp2021}, toda sesi\'on de elicitaci\'on requiere consentimiento libre, espec\'ifico, informado e inequ\'ivoco del participante. Ning\'un dato personal de terceros (n\'umeros de c\'edula, tel\'efonos, correos personales, grabaciones de audio o actas firmadas) se almacena en el repositorio p\'ublico de GitHub; dichas evidencias se resguardan en la carpeta institucional privada (\texttt{@uteq.edu.ec}) compartida \'unicamente con el docente, y en el repositorio se publicar\'a el manifiesto con la huella criptogr\'afica SHA-256 correspondiente.

\begin{tcolorbox}[colback=boxbg,colframe=midblue,boxrule=0.8pt,arc=2mm,breakable,title={\textbf{Plantilla de Consentimiento Informado para Registro de Audio (ACTA-CONS-01)}}]
\textbf{Proyecto:} Plataforma de Ventas con M\'odulo de Recomendaci\'on basado en IA para Establecimientos de Suministros El\'ectricos en Quevedo.\\
\textbf{Asignatura:} Ingenier\'ia de Requerimientos (ISR-401) --- Universidad T\'ecnica Estatal de Quevedo.\\[4pt]
\textbf{1. Finalidad del registro:} Comprender el flujo diario de trabajo en mostrador, los tiempos de atenci\'on y los criterios de asesor\'ia t\'ecnica con fines estrictamente acad\'emicos, sin evaluar el desempe\~no laboral individual.\\[3pt]
\textbf{2. Voluntariedad y revocaci\'on:} La participaci\'on es libre y voluntaria. El participante puede omitir cualquier pregunta, solicitar una pausa o retirar su consentimiento en cualquier momento sin consecuencia alguna.\\[3pt]
\textbf{3. Anonimizaci\'on y custodia:} Los nombres reales del participante y del local comercial se reemplazan por c\'odigos de rol (por ejemplo, \texttt{SH-01}). El archivo de audio se conserva exclusivamente en almacenamiento institucional privado UTEQ bajo los principios de confidencialidad y minimizaci\'on de datos de la LOPDP \cite{lopdp2021}.\\[6pt]
\noindent\begin{tabularx}{\linewidth}{|Y|Y|}
\hline
\textbf{C\'odigo de sesi\'on:} \texttt{GA-TL-02-ENS-01} & \textbf{Fecha:} \_\_\_ / \_\_\_ / 2026\\
\hline
\textbf{Rol representado:} Cajero / Vendedor (SH-01) & \textbf{Autoriza grabaci\'on de audio:} [ \ ] S\'i \quad [ \ ] No\\
\hline
\textbf{Nombre del entrevistador:} \rule{0pt}{14pt} & \textbf{Firma / confirmaci\'on verbal:} \rule{0pt}{14pt}\\
\hline
\end{tabularx}
\end{tcolorbox}


% ==============================================================================
% SECCION 2: GUION DE ENTREVISTA ESTRUCTURADA (Integrante 1)
% ==============================================================================
\section{Guion de entrevista estructurada (SH-01: Cajero y Vendedor de Mostrador)}

\subsection{Evoluci\'on metodol\'ogica del instrumento}
El presente guion deriva de la consolidaci\'on entre la versi\'on preliminar (\emph{Gu\'ia de preguntas para entrevista v3}) y el instrumento formal de campo (\emph{Gu\'ia Definitiva de Entrevistas y Observaci\'on en Sitio, ACTA-ELI-01}). Para atender los criterios de elicitaci\'on de la norma ISO/IEC/IEEE 29148:2018 \cite{iso29148} y las perspectivas de contexto de Pohl \cite{pohl2025}, se aplicaron tres mejoras concretas sobre la versi\'on v3:

\begin{enumerate}[leftmargin=18pt,itemsep=2pt]
    \item \textbf{Segmentaci\'on por parte interesada:} La versi\'on v3 agrupaba 29 preguntas en una lista continua que mezclaba tareas de mostrador con decisiones administrativas (pedidos a proveedores, estacionalidad de compras y pol\'iticas de stock). Se aislaron \'unicamente las preguntas pertinentes al perfil \textbf{SH-01 (Cajero / Vendedor de Mostrador)}.
    \item \textbf{Incorporaci\'on del contexto operativo de Quevedo:} Se a\~nadieron preguntas sobre ergonom\'ia f\'isica del mostrador (uso de teclado y lector frente al rat\'on), jerga el\'ectrica coloquial local (\emph{breaker}, \emph{taipe}, \emph{boquilla}, \emph{contactora}, \emph{manguera}) y contingencias por intermitencia de internet o ca\'idas del SRI.
    \item \textbf{Formalizaci\'on de estructura y acotaci\'on:} Se integraron el protocolo real de apertura con consentimiento informado y el cierre estructurado del artefacto \texttt{ACTA-ELI-01}, sintetizando el cuestionario en \textbf{15 preguntas estructuradas} organizadas en cuatro bloques tem\'aticos para el ensayo controlado de \texttt{GA-TL-02}.
\end{enumerate}

\subsection{Ficha de control del instrumento}
\noindent\begin{tabularx}{\linewidth}{|p{0.28\linewidth}|Y|p{0.24\linewidth}|Y|}
\hline
\thdr{C\'odigo de acta} & \texttt{ACTA-ELI-CAJ-01} & \thdr{Duraci\'on estimada} & 30 a 45 minutos\\
\hline
\textbf{Parte interesada} & Cajero / Vendedor (SH-01) & \textbf{Tipo de entrevista} & Estructurada (15 preguntas)\\
\hline
\rowcolor{rowgray} \textbf{Nivel electrot\'ecnico} & [ \ ] Emp\'irico [ \ ] T\'ecnico [ \ ] Ninguno & \textbf{Autoriza audio} & [ \ ] S\'i \quad [ \ ] No\\
\hline
\end{tabularx}

\subsection{Guion de apertura y consentimiento verbal}
\begin{tcolorbox}[colback=boxbg,colframe=uteqgreen,boxrule=0.7pt,arc=1.5mm,breakable]
\emph{``Buenos d\'ias/tardes. Mi nombre es Gilces Carranza Jos\'e Ignacio, estudiante de Ingenier\'ia de Software de la Universidad T\'ecnica Estatal de Quevedo. Estamos desarrollando un proyecto acad\'emico sobre una Plataforma de Ventas con asistencia inteligente para negocios de materiales el\'ectricos en Quevedo. El objetivo es comprender c\'omo trabajan actualmente en su d\'ia a d\'ia, qu\'e problemas enfrentan y qu\'e caracter\'isticas har\'ian realmente \'util una herramienta as\'i. Su participaci\'on es voluntaria, puede no responder alguna pregunta o terminar la entrevista cuando lo desee. No hay respuestas correctas o incorrectas, y la informaci\'on se usar\'a estrictamente con fines acad\'emicos manteniendo el anonimato del local y del personal. ¿Acepta participar y me permite grabar la conversaci\'on para no perder detalles t\'ecnicos?''}
\end{tcolorbox}

\subsection{Bloques tem\'aticos de preguntas (15 preguntas)}

\subsubsection*{Bloque 1: Perfil operativo, rutina y ergonom\'ia en mostrador (Perspectiva de uso y f\'isica)}
\begin{enumerate}[label=\textbf{P\arabic*.},leftmargin=26pt,itemsep=3pt]
    \item ¿Cu\'ales son sus responsabilidades exactas durante el turno (cobrar, despachar, acomodar perchas o revisar inventario) desde que abre hasta que cierra?\\
    {\footnotesize\textit{Objetivo anal\'itico:} Medir la dispersi\'on de atenci\'on del cajero y las interrupciones mientras opera en caja.}
    \item Descr\'ibame paso a paso qu\'e ocurre desde que un cliente llega al mostrador hasta que se retira con su compra y comprobante.\\
    {\footnotesize\textit{Objetivo anal\'itico:} Mapear el flujo actual (\emph{As-Is}) e identificar tiempos muertos entre atenci\'on, b\'usqueda en bodega y cobro.}
    \item ¿En qu\'e horarios del d\'ia se concentran las mayores filas, cu\'antos clientes llegan a esperar y qu\'e pasos le toman m\'as tiempo en horas pico?\\
    {\footnotesize\textit{Objetivo anal\'itico:} Identificar picos de estr\'es operativo y cuellos de botella cr\'iticos.}
    \item ¿C\'omo localiza un producto en el sistema y en bodega, y qu\'e tan frecuente es que el sistema indique existencias pero f\'isicamente est\'e agotado?\\
    {\footnotesize\textit{Objetivo anal\'itico:} Verificar el uso de perif\'ericos (lector de c\'odigo de barras, teclado), ergonom\'ia del mostrador y desfases de inventario.}
\end{enumerate}

\subsubsection*{Bloque 2: Herramienta actual, lenguaje local y perfil de clientes (Perspectiva de sociedad y lenguaje)}
\begin{enumerate}[label=\textbf{P\arabic*.},resume,leftmargin=26pt,itemsep=3pt]
    \item ¿Qu\'e herramienta usan hoy para registrar ventas y emitir facturas, qu\'e le genera mayor fatiga durante el cobro y qu\'e pas\'o la \'ultima vez que una venta se complic\'o?\\
    {\footnotesize\textit{Objetivo anal\'itico:} Detectar deficiencias de usabilidad e incidentes cr\'iticos por errores de digitaci\'on o lentitud.}
    \item ¿Qui\'enes son sus clientes m\'as frecuentes (maestros electricistas, contratistas, particulares, empresas) y c\'omo cambia su forma de atender a cada perfil?\\
    {\footnotesize\textit{Objetivo anal\'itico:} Calibrar el nivel t\'ecnico requerido en la atenci\'on seg\'un el tipo de comprador.}
    \item ¿Con qu\'e t\'erminos coloquiales o ``nombres de calle'' piden los materiales en Quevedo (por ejemplo, \emph{breaker}, \emph{taipe}, \emph{boquilla}, \emph{contactora}, \emph{manguera}) y cu\'ales generan mayor confusi\'on?\\
    {\footnotesize\textit{Objetivo anal\'itico:} Alimentar el diccionario de sin\'onimos del dominio para el buscador y el m\'odulo de recomendaci\'on.}
\end{enumerate}

\subsubsection*{Bloque 3: Heur\'istica de asesor\'ia t\'ecnica, sustitutos y venta cruzada (L\'ogica del m\'odulo IA)}
\begin{enumerate}[label=\textbf{P\arabic*.},resume,leftmargin=26pt,itemsep=3pt]
    \item ¿Hay clientes que buscan siempre a un vendedor espec\'ifico para que los asesore y por qu\'e cree que ocurre eso?\\
    {\footnotesize\textit{Objetivo anal\'itico:} Detectar conocimiento t\'acito concentrado en vendedores con mayor experiencia.}
    \item ¿Con qu\'e frecuencia los clientes le piden consejo sobre calibres de cable, amperajes de \emph{breaker}, voltajes o compatibilidad antes de comprar?\\
    {\footnotesize\textit{Objetivo anal\'itico:} Cuantificar la demanda real de asesor\'ia t\'ecnica en mostrador.}
    \item Cuando un cliente solicita una recomendaci\'on para una instalaci\'on residencial o industrial, ¿qu\'e datos le pide usted antes de sugerirle un producto?\\
    {\footnotesize\textit{Objetivo anal\'itico:} Elicitar las variables de diagn\'ostico de entrada para el asistente de recomendaci\'on.}
    \item Cuando el producto solicitado est\'a agotado, ¿en qu\'e criterios t\'ecnicos se basa para ofrecer un sustituto sin arriesgarse a provocar un cortocircuito o un reclamo?\\
    {\footnotesize\textit{Objetivo anal\'itico:} Extraer reglas de equivalencia t\'ecnica (tensi\'on, corriente, curva, dimensiones y marca).}
    \item ¿Suele ofrecer productos complementarios (cinta aislante, cajas t\'ermicas, terminales) y c\'omo sabe cu\'ales sugerir en cada compra?\\
    {\footnotesize\textit{Objetivo anal\'itico:} Definir reglas de asociaci\'on para venta cruzada y reducci\'on de compras incompletas.}
\end{enumerate}

\subsubsection*{Bloque 4: Aprendizaje, interfaz de apoyo y contingencias de red (Perspectiva t\'ecnica y TI)}
\begin{enumerate}[label=\textbf{P\arabic*.},resume,leftmargin=26pt,itemsep=3pt]
    \item ¿Cu\'anto tiempo le toma a un vendedor nuevo dominar las equivalencias del cat\'alogo y a qui\'en o a qu\'e recurre usted cuando tiene dudas de compatibilidad?\\
    {\footnotesize\textit{Objetivo anal\'itico:} Medir la curva de aprendizaje y el tiempo invertido en consultas manuales.}
    \item Si el sistema mostrara sugerencias de sustitutos o complementos al vender, ¿en qu\'e momento le resultar\'ia \'util verlas y qu\'e datos necesita en pantalla para confiar en esa recomendaci\'on?\\
    {\footnotesize\textit{Objetivo anal\'itico:} Definir el punto de interacci\'on en caja y los atributos de explicabilidad (precio, stock, especificaciones y motivo).}
    \item ¿Con qu\'e frecuencia falla el internet o se satura el servicio del SRI durante la semana y qu\'e procedimiento siguen en ese instante para no detener el cobro?\\
    {\footnotesize\textit{Objetivo anal\'itico:} Fundamentar los requisitos de contingencia, operaci\'on sin conexi\'on y degradaci\'on controlada.}
\end{enumerate}

\subsection{Guion de cierre}
\begin{tcolorbox}[colback=boxbg,colframe=uteqgreen,boxrule=0.7pt,arc=1.5mm,breakable]
\emph{``Muchas gracias por su tiempo y por compartir su experiencia con nosotros. Toda esta informaci\'on nos ayudar\'a a dise\~nar una soluci\'on ajustada a la realidad de los locales el\'ectricos de Quevedo. ¿Nos permitir\'ia contactarle m\'as adelante para mostrarle un prototipo inicial y validar su funcionamiento?''}
\end{tcolorbox}


% ==============================================================================
% SECCION 3: MAPA DE EMPATIA Y DEFINICION DEL PROBLEMA (Integrante 2)
% ==============================================================================
\section{Design Thinking: Mapa de Empat\'ia y definici\'on del problema}

\subsection{Prop\'osito del lienzo}
Tomando el \emph{design thinking} como base para la ingenier\'ia de requisitos \cite{ireb2026}, en esta secci\'on se arma el Mapa de Empat\'ia del \textbf{Cajero / Vendedor de Mostrador (SH-01)}. Se usa para ordenar lo que salga del ensayo de entrevista con ese rol: qu\'e dice, qu\'e hace, qu\'e piensa y qu\'e siente cuando atiende en el mostrador. Las celdas se llenaron con lo que se dijo en el ensayo \texttt{GA-TL-02-ENS-01}, cuyo audio est\'a registrado en \texttt{GA/TL-02/evidencias/MANIFIESTO.md}. Las frases entre comillas son textuales.

\subsection{Mapa de Empat\'ia (Perfil: SH-01 --- Cajero / Vendedor de Mostrador)}

\begin{tcolorbox}[colback=boxbg,colframe=midblue,boxrule=0.8pt,arc=2mm,title={\textbf{Perfil analizado: Cajero / Vendedor de Mostrador (SH-01) --- Ferreter\'ia el\'ectrica en Quevedo}}]
\textbf{Escenario ensayado:} Atenci\'on en mostrador durante horas pico (ma\~nana y mediod\'ia), b\'usqueda de materiales solicitados con jerga local, asesor\'ia t\'ecnica sobre reemplazos y cobro con facturaci\'on electr\'onica.\\
\textbf{Fuente de respaldo:} Audio del ensayo \texttt{GA-TL-02-ENS-01} registrado en \texttt{GA/TL-02/evidencias/MANIFIESTO.md}.
\end{tcolorbox}

\vspace{4pt}
{\small
\noindent\begin{tabularx}{\linewidth}{|Y|Y|}
\hline
\thdr{¿Qu\'e DICE? (Expresiones del ensayo)} & \thdr{¿Qu\'e HACE? (Acciones en mostrador)}\\
\hline
\begin{itemize}[leftmargin=12pt,itemsep=2pt,topsep=2pt]
    \item ``La gente llega pidiendo con nombres de calle como \emph{taipe}, \emph{breaker} o \emph{boquilla}, y a veces toca adivinar el amperaje.'' \textit{(ENS-01)}
    \item ``Cuando se hace fila al mediod\'ia, perder tres minutos buscando un reemplazo en bodega o preguntando al due\~no molesta al cliente.'' \textit{(ENS-01)}
    \item ``Si el internet se cae o el SRI se pone lento, toca anotar en libreta y cuadrar despu\'es.'' \textit{(ENS-01)}
\end{itemize}
&
\begin{itemize}[leftmargin=12pt,itemsep=2pt,topsep=2pt]
    \item Alterna constantemente entre cobrar en caja, revisar existencias en bodega y responder dudas t\'ecnicas de clientes particulares.
    \item Consulta al administrador, revisa cat\'alogos impresos o busca en el tel\'efono cuando desconoce una equivalencia el\'ectrica.
    \item Prioriza el uso de teclado y lector de barras por el espacio reducido en el mostrador.
\end{itemize}\\
\hline
\thdr{¿Qu\'e PIENSA Y SIENTE? (Emociones y temores)} & \thdr{¿Qu\'e VE y OYE? (Entorno f\'isico y social)}\\
\hline
\begin{itemize}[leftmargin=12pt,itemsep=2pt,topsep=2pt]
    \item Siente presi\'on cuando se acumulan varios clientes apurados en la fila y debe verificar datos t\'ecnicos al mismo tiempo.
    \item Teme recomendar un material incompatible que provoque una devoluci\'on, un da\~no en la instalaci\'on del cliente o un reclamo del administrador.
    \item Percibe que un vendedor nuevo tarda meses en memorizar equivalencias y nombres locales del cat\'alogo.
\end{itemize}
&
\begin{itemize}[leftmargin=12pt,itemsep=2pt,topsep=2pt]
    \item Observa un mostrador estrecho, cajas de materiales el\'ectricos y clientes revisando la hora mientras esperan su factura.
    \item Escucha pedidos apurados de maestros electricistas que exigen precisi\'on y clientes particulares que no saben el nombre t\'ecnico del producto.
    \item Oye quejas cuando un producto aparece con saldo en pantalla pero no existe f\'isicamente en percha.
\end{itemize}\\
\hline
\thdr{DOLORES / FRUSTRACIONES} & \thdr{RESULTADOS / NECESIDADES}\\
\hline
\begin{itemize}[leftmargin=12pt,itemsep=2pt,topsep=2pt]
    \item Demoras de varios minutos para confirmar compatibilidad de sustitutos cuando un art\'iculo est\'a agotado.
    \item Desfases entre el inventario registrado y la disponibilidad real en bodega.
    \item Interrupciones operativas cuando falla la conectividad o el servicio tributario del SRI.
\end{itemize}
&
\begin{itemize}[leftmargin=12pt,itemsep=2pt,topsep=2pt]
    \item Responder con rapidez y seguridad sobre alternativas compatibles sin depender siempre del vendedor m\'as antiguo.
    \item Confiar en que el stock consultado corresponde a la existencia f\'isica real.
    \item Mantener la continuidad de la atenci\'on en caja incluso ante fallos externos de red.
\end{itemize}\\
\hline
\end{tabularx}
}

\subsection{Definici\'on del problema centrada en el usuario}
La definici\'on del problema se redacta estrictamente desde la necesidad operativa del trabajador en mostrador, evitando imponer una soluci\'on tecnol\'ogica espec\'ifica:

\begin{tcolorbox}[colback=white,colframe=uteqgreen,boxrule=0.8pt,arc=1.5mm]
\textbf{Enunciado de punto de vista (POV):}\\
El cajero y vendedor de mostrador de un establecimiento de suministros el\'ectricos en Quevedo necesita identificar en segundos productos compatibles, disponibles en bodega y reconocibles por su nombre coloquial local, porque durante las horas pico las consultas t\'ecnicas manuales y los desfases de inventario retrasan la atenci\'on y generan riesgo de vender materiales incompatibles.\\[8pt]
\textbf{Pregunta generadora («¿C\'omo podr\'iamos...?»):}\\
\emph{¿C\'omo podr\'iamos ayudar al cajero y vendedor de mostrador a identificar con rapidez y certeza alternativas el\'ectricas compatibles y disponibles —incluso cuando el cliente usa nombres coloquiales— sin detener el flujo de atenci\'on en horas pico ni depender de consultas manuales prolongadas?}
\end{tcolorbox}


% ==============================================================================
% SECCIONES PROGRAMADAS PARA EL ENSAYO Y CIERRE (Dias 2 y 3)
% ==============================================================================
% \section{An\'alisis y correcci\'on de preguntas defectuosas}

\subsection{Criterios de auditor\'ia tras el ensayo}
Durante el ensayo de juego de roles se detectaron interrupciones en la fluidez de varias respuestas: el participante omiti\'o partes de las preguntas compuestas, respondi\'o con monos\'ilabos ante formulaciones cerradas o repiti\'o el riesgo sugerido por el entrevistador. Con base en las recomendaciones de elicitaci\'on de Pohl \cite{pohl2025}, el temario IREB \cite{ireb2026} y los criterios de no ambig\"uedad de ISO/IEC/IEEE 29148:2018 \cite{iso29148}, se auditaron las 15 preguntas originales y se identificaron \textbf{ocho preguntas defectuosas} clasificadas en cuatro categor\'ias: compuestas (dobles o triples), cerradas sin necesidad, inductivas y centradas en la soluci\'on.

\subsection{Cat\'alogo de preguntas que fallaron y su versi\'on corregida}

{\footnotesize
\noindent\begin{tabularx}{\linewidth}{|p{0.06\linewidth}|p{0.27\linewidth}|p{0.17\linewidth}|Y|}
\hline
\thdr{ID} & \thdr{Pregunta original (fall\'o en ensayo)} & \thdr{Defecto detectado} & \thdr{Versi\'on corregida (abierta, neutral y simple)}\\
\hline
\textbf{P3} & ?`En qu\'e horarios del d\'ia se concentran las mayores filas, cu\'antos clientes llegan a esperar y qu\'e pasos le toman m\'as tiempo en horas pico? & \textbf{Pregunta triple:} solicita horarios, volumen de espera y cuellos de botella en un solo enunciado. & \textbf{P3a:} ?`En qu\'e franjas horarias se acumula mayor cantidad de clientes en el mostrador? \newline \textbf{P3b:} Durante esas horas pico, ?`qu\'e pasos de la atenci\'on consumen m\'as tiempo?\\
\hline
\rowcolor{rowgray} \textbf{P4} & ?`C\'omo localiza un producto en el sistema y en bodega, y qu\'e tan frecuente es que el sistema indique existencias pero f\'isicamente est\'e agotado? & \textbf{Pregunta doble:} mezcla el mecanismo de b\'usqueda con la frecuencia de desfases de inventario. & \textbf{P4a:} ?`Qu\'e pasos sigue para localizar un producto en el sistema y ubicarlo f\'isicamente en bodega? \newline \textbf{P4b:} ?`Qu\'e ocurre cuando el registro indica disponibilidad pero el producto no aparece en percha?\\
\hline
\textbf{P5} & ?`Qu\'e herramienta usan hoy para registrar ventas y emitir facturas, qu\'e le genera mayor fatiga durante el cobro y qu\'e pas\'o la \'ultima vez que una venta se complic\'o? & \textbf{Pregunta triple:} combina herramientas actuales, fatiga operativa y relato de un incidente cr\'itico. & \textbf{P5a:} ?`Qu\'e herramientas utiliza actualmente para registrar ventas y emitir comprobantes? \newline \textbf{P5b:} Cu\'enteme qu\'e ocurri\'o la \'ultima vez que una venta se retras\'o o se complic\'o en caja.\\
\hline
\rowcolor{rowgray} \textbf{P8} & ?`Hay clientes que buscan siempre a un vendedor espec\'ifico para que los asesore y por qu\'e cree que ocurre eso? & \textbf{Cerrada al inicio y doble:} induce una respuesta de s\'i/no antes de explorar la causa. & ?`De qu\'e manera influye la experiencia de cada vendedor cuando un cliente solicita asesor\'ia t\'ecnica en el mostrador?\\
\hline
\textbf{P11} & Cuando el producto solicitado est\'a agotado, ?`en qu\'e criterios t\'ecnicos se basa para ofrecer un sustituto sin arriesgarse a provocar un cortocircuito o un reclamo? & \textbf{Inductiva:} condiciona la respuesta mencionando expl\'icitamente \enquote{cortocircuito} y \enquote{reclamo}. & Cuando un producto solicitado est\'a agotado, ?`qu\'e criterios eval\'ua usted para decidir qu\'e alternativa ofrecer al cliente?\\
\hline
\rowcolor{rowgray} \textbf{P12} & ?`Suele ofrecer productos complementarios (cinta aislante, cajas t\'ermicas, terminales) y c\'omo sabe cu\'ales sugerir en cada compra? & \textbf{Cerrada, doble y sesgada:} enumera ejemplos concretos que orientan al entrevistado. & ?`C\'omo decide qu\'e productos adicionales o complementarios sugerir junto con el material principal que pide el cliente?\\
\hline
\textbf{P14} & Si el sistema mostrara sugerencias de sustitutos o complementos al vender, ?`en qu\'e momento le resultar\'ia \'util verlas y qu\'e datos necesita en pantalla para confiar en esa recomendaci\'on? & \textbf{Centrada en la soluci\'on y doble:} asume una interfaz hipot\'etica antes de validar la necesidad real. & \textbf{P14a:} Al momento de atender en mostrador, ?`qu\'e informaci\'on t\'ecnica necesita verificar para estar seguro de que un reemplazo es compatible? \newline \textbf{P14b:} ?`En qu\'e instante de la venta le resulta necesario consultar esa informaci\'on?\\
\hline
\rowcolor{rowgray} \textbf{P15} & ?`Con qu\'e frecuencia falla el internet o se satura el servicio del SRI durante la semana y qu\'e procedimiento siguen en ese instante para no detener el cobro? & \textbf{Pregunta doble:} une frecuencia de fallos externos con protocolo de contingencia operativa. & \textbf{P15a:} ?`Con qu\'e frecuencia experimentan interrupciones de internet o lentitud en el SRI durante la semana? \newline \textbf{P15b:} Cuando ocurre una interrupci\'on de ese tipo, ?`qu\'e procedimiento siguen en caja para continuar atendiendo?\\
\hline
\end{tabularx}
}

\subsection{Ajustes aplicados para la entrevista real de campo (PE-U2 y PFC)}
La auditor\'ia demostr\'o que agrupar varias consultas en un solo enunciado provoca que el participante omita detalles relevantes sobre los tiempos de espera y los procedimientos de contingencia. Para las sesiones reales en los establecimientos de suministros el\'ectricos de Quevedo se aplicar\'an preguntas at\'omicas de una sola idea, se retirar\'an t\'erminos inductivos y se usar\'an repreguntas breves (\emph{\enquote{?`me puede dar un ejemplo reciente?}}) solo cuando se requiera ampliar una respuesta.

% \section{Manifiesto de evidencias y declaraci\'on de uso de IA}

\subsection{Manifiesto del audio del ensayo}

La evidencia del roleplay se identifica con el c\'odigo de sesi\'on \texttt{GA-TL-02-ENS-01}
y corresponde al perfil \textbf{SH-01 (Cajero / Vendedor de Mostrador)}. Sus datos se
registran en el repositorio sin publicar el archivo de audio.

\noindent\begin{tabularx}{\linewidth}{|p{0.32\linewidth}|Y|}
\hline
\thdr{Campo} & \thdr{Dato verificado}\\
\hline
Nombre del archivo & \texttt{ensayo\_entrevista\_SH01\_20261120.mp3}\\
\hline
\rowcolor{rowgray} Fecha de grabaci\'on & 20 de noviembre de 2026\\
\hline
Duraci\'on & 10:35\\
\hline
\rowcolor{rowgray} Formato & MP3\\
\hline
SHA-256 & \texttt{\footnotesize 31986071108A8BD445B3DCF5C045BB5EF207288A36967A2E2DDD1CAF397FDDA1}\\
\hline
\end{tabularx}

\vspace{4pt}
La huella se obtuvo con \texttt{Get-FileHash \ldots -Algorithm SHA256} sobre el archivo
final, sin conversiones posteriores. El audio se inicia con el guion de apertura y la
autorizaci\'on verbal del participante, seg\'un el protocolo de consentimiento descrito en
la secci\'on de contexto y en el registro \texttt{ACTA-CONS-01}.

\subsection{Custodia de la evidencia}

De acuerdo con la LOPDP del Ecuador \cite{lopdp2021}, el archivo de audio se conserva en
una carpeta privada de OneDrive institucional (\texttt{@uteq.edu.ec}) del integrante
responsable, compartida con el docente \texttt{gguerrero@uteq.edu.ec}. No se publica en el
repositorio ning\'un dato personal de terceros: ni c\'edulas, tel\'efonos, correos personales,
audio ni consentimientos firmados. El registro detallado se encuentra en
\texttt{GA/TL-02/evidencias/MANIFIESTO.md}.

\subsection{Declaraci\'on de uso de IA como apoyo}

\begin{itemize}
    \item Herramienta: OpenCode, modelo \texttt{Big Pickle} (proveedor \texttt{opencode}).
    \item Versi\'on: modelo \texttt{Big Pickle}, sesi\'on de los d\'ias 20 y 21 de noviembre de 2026.
    \item Prop\'osito: apoyo en la estructura \LaTeX{}, en la revisi\'on de las normas de entrega
          y en el c\'alculo de la huella SHA-256 del audio del ensayo.
    \item Verificaci\'on: el equipo revis\'o manualmente cada secci\'on, contrast\'o los datos del
          manifiesto con el archivo original mediante \texttt{Get-FileHash} y valid\'o la
          compilaci\'on local con MiKTeX. Las respuestas y apreciaciones del roleplay provienen
          de la grabaci\'on real del equipo.
\end{itemize}


\newpage
\printbibliography[title={Referencias}]

\end{document}
